%!TEX root = ../thesis.tex
%*******************************************************************************
%****************************** Second Chapter *********************************
%*******************************************************************************

\chapter{Methodology}

\ifpdf
    \graphicspath{{Chapter2/Figs/Raster/}{Chapter2/Figs/PDF/}{Chapter2/Figs/}}
\else
    \graphicspath{{Chapter2/Figs/Vector/}{Chapter2/Figs/}}
\fi

The availability of different datasets and sensing technologies enables to rethink the role of data in civic participation. This provides opportunities to focus on digital technology and mechanisms that support human centred and alternative approaches. Taking a Participatory Action Research (PAR) approach as a co-design processes and applying iterative design and more engineering focused design sprints, this work undertakes the following process to study the use of data and technology in civic participation and advocacy:

\begin{itemize}
  \item Identifying the need for data (Discovery Phase). Through conducting interviews and observations of existing groups map out how they are using/not using data. 
  \item Making data and technology accessible (Design/Develop Phase). Undertaking a co-design process to make platforms and tools that respond to the needs of citizens. Starting with paper based methods and then moving to digital.
  \item Meaningful presentation of data (Deploy Phase). Ascertain what visualisations and formats of data are more accessible and useful for the sensemaking processes in different contexts. 
  \item Taking collective action by purposing data (Sensemaking Phase). Supporting the processes of sensemaking and assessing the use of data for decision-making and developing new solutions.
  \item Assessing the outcomes of data use (Evaluation Phase). Evaluating the use of data for improving civic participation and turning it to activities for local benefit.
\end{itemize}

%FIGURE framework HERE
\begin{figure}
\centering
 \includegraphics[scale=0.4]{framework}
  \caption{Illustration of the the four steps of the toolkit}
  \label{fig:framework}
\end{figure}
